\documentclass[12pt, a4paper]{article}

% Pakiety
\usepackage[utf8]{inputenc}
\usepackage[T1]{fontenc}
\usepackage[polish]{babel}
\usepackage{amsmath, amssymb}
\usepackage{graphicx}
\usepackage{float}
\usepackage{geometry}
\usepackage{hyperref}
\usepackage{booktabs}
\usepackage{caption}
\usepackage{listings}
\usepackage{xcolor}
\usepackage{parskip}

\geometry{
  a4paper,
  left=2.5cm,
  right=2.5cm,
  top=2.5cm,
  bottom=2.5cm
}

% Zakaz dzielenia wyrazów
\hyphenpenalty=10000
\exhyphenpenalty=10000
\emergencystretch=3em

\hypersetup{
  colorlinks=true,
  linkcolor=blue,
  urlcolor=blue,
  citecolor=blue
}

% Zmiana nazwy float figure: Rysunek -> Wykres
\addto\captionspolish{\renewcommand{\figurename}{Wykres}}

% Styl listingów
\lstset{
  language=Python,
  basicstyle=\ttfamily\small,
  keywordstyle=\color{blue}\bfseries,
  commentstyle=\color{gray}\itshape,
  stringstyle=\color{orange!80!black},
  numbers=left,
  numberstyle=\tiny\color{gray},
  numbersep=8pt,
  frame=single,
  breaklines=true,
  breakatwhitespace=true,
  captionpos=b,
  showstringspaces=false,
  columns=flexible,
  keepspaces=true,
  inputencoding=utf8,
  extendedchars=true,
  literate={ą}{{\k{a}}}1 {ę}{{\k{e}}}1 {ó}{{\'{o}}}1
            {ś}{{\'{s}}}1 {ł}{{\l{}}}1 {ź}{{\'{z}}}1
            {ż}{{\.z}}1 {ć}{{\'{c}}}1 {ń}{{\'{n}}}1
            {Ą}{{\k{A}}}1 {Ę}{{\k{E}}}1 {Ó}{{\'{O}}}1
            {Ś}{{\'{S}}}1 {Ł}{{L}}1,
}

% Strona tytułowa
\title{
  \textbf{Metody Obliczeniowe w Nauce i Technice} \\[0.5em]
  \large Laboratorium 11: Optymalizacja
}
\author{Mateusz Stoch}
\date{27 maja 2026}

\begin{document}

\maketitle

% ============================================================
\section{Zadanie 1: Preconditioning}
% ============================================================

\subsection{Zadanie 1(a): Gradient prosty z optymalnym krokiem}

\textbf{Treść zadania:} Znajdź minimum funkcji $f$ metodą gradientu prostego, wyznaczając w każdej iteracji optymalną wartość kroku $\alpha_k$. Narysuj wykres konturowy funkcji $f$, na którym zaznacz iteracje algorytmu. Jako punkt startowy przyjmij $(x_0, y_0) = (9, 1)$.

Funkcja celu jest określona wzorem:
\[
  f(x, y) = \frac{1}{2}x^2 + \frac{9}{2}y^2
\]

Optymalną długość kroku w każdej iteracji możemy wyznaczyć analitycznie za pomocą dokładnego przeszukiwania kierunkowego (exact line search). Wzór ten wyprowadza się poprzez znalezienie $\alpha$, które minimalizuje wartość funkcji w nowym punkcie, czyli szukamy minimum funkcji jednej zmiennej $h(\alpha) = f(x_k - \alpha g_k)$. Po zróżniczkowaniu tej funkcji względem $\alpha$ i przyrównaniu pochodnej do zera otrzymujemy optymalny krok $\alpha_k$ dla funkcji kwadratowej o Hesjanie $H$:
\[
  \alpha_k = \frac{\|g_k\|^2}{g_k^T H g_k}
\]
gdzie $g_k = \nabla f(x_k, y_k) = [x_k, 9y_k]^T$ to gradient w danej iteracji, a $H = \nabla^2 f(x, y) = \begin{bmatrix} 1 & 0 \\ 0 & 9 \end{bmatrix}$ to stała macierz drugich pochodnych (Hesjan).

\begin{lstlisting}[caption={Gradient prosty z optymalnym krokiem}]
def optimal_step_f1(xy):
    g = grad_f1(xy)
    H = np.diag([1.0, 9.0])
    return (g @ g) / (g @ H @ g)

def gradient_descent_f1(xy0, tol=1e-10, max_iter=1000):
    xy = np.array(xy0, dtype=float)
    history = [xy.copy()]
    for _ in range(max_iter):
        g = grad_f1(xy)
        if np.linalg.norm(g) < tol:
            break
        alpha = optimal_step_f1(xy)
        xy = xy - alpha * g
        history.append(xy.copy())
    return np.array(history)
\end{lstlisting}

Uruchomienie algorytmu z punktu startowego $(9, 1)$ dało następujące wyniki:
\begin{itemize}
  \item Liczba iteracji potrzebna do osiągnięcia zbieżności (z tolerancją $\|g_k\| < 10^{-10}$): 115
  \item Znalezione minimum: $x = 0{,}000000$, $y = 0{,}000000$
  \item Wartość funkcji w minimum: $2{,}31 \cdot 10^{-21}$
\end{itemize}

Trajektorię punktów pośrednich przedstawiono na Wykresie~\ref{fig:wykres1}. Charakterystyczny zygzakowaty kształt wynika z faktu, że kierunek najszybszego spadku (gradient) nie jest skierowany bezpośrednio na minimum, lecz jest prostopadły do linii poziomicowych (konturów), które są mocno wydłużonymi elipsami.

\begin{figure}[H]
  \centering
  \includegraphics[width=0.85\textwidth]{wykres1.png}
  \caption{Wykres konturowy funkcji $f(x,y)$ wraz z zaznaczoną zygzakowatą trajektorią metody gradientu prostego.}
  \label{fig:wykres1}
\end{figure}

\subsection{Zadanie 1(b): Uwarunkowanie Hessiana}

\textbf{Treść zadania:} Oblicz współczynnik uwarunkowania macierzy $\nabla^2 f(x, y)$.

Hesjan funkcji $f(x, y)$ jest stały i wynosi:
\[
  \nabla^2 f(x, y) = H = \begin{bmatrix} 1 & 0 \\ 0 & 9 \end{bmatrix}
\]

Wartości własne tej macierzy diagonalnej to bezpośrednio jej elementy na przekątnej:
\[
  \lambda_{\min} = 1{,}0, \qquad \lambda_{\max} = 9{,}0
\]

Współczynnik uwarunkowania macierzy $\kappa$ definiuje się jako stosunek największej do najmniejszej wartości własnej (dla macierzy symetrycznych dodatnio określonych):
\[
  \kappa = \frac{\lambda_{\max}}{\lambda_{\min}} = \frac{9{,}0}{1{,}0} = 9{,}0
\]

Współczynnik ten mierzy, jak bardzo elipsy poziomicowe są rozciągnięte. Im większy jest $\kappa$, tym bardziej wydłużone elipsy, co spowalnia zbieżność gradientu prostego. Teoretyczny mnożnik zbieżności (ogólne oszacowanie tempa spadku błędu na krok, szczegółowo omówione w pracy J. R. Shewchuka~\cite{shewchuk_cg}) wynosi:
\[
  R = \left( \frac{\kappa - 1}{\kappa + 1} \right)^2 = \left( \frac{9 - 1}{9 + 1} \right)^2 = \left( \frac{8}{10} \right)^2 = 0{,}64
\]
Oznacza to, że w każdym kroku błąd może zmniejszyć się co najwyżej o czynnik $0{,}64$, co tłumaczy dużą liczbę iteracji (115) potrzebną do osiągnięcia bardzo wysokiej dokładności.

\subsection{Zadanie 1(c): Dekompozycja Choleskiego Hessiana}

\textbf{Treść zadania:} Znajdź macierz $L$ taką, że: $\nabla^2 f(x, y) = L L^T$.

Poszukujemy macierzy dolnotrójkątnej $L$, która po przemnożeniu przez swoją transpozycję da macierz Hessiana. Stosując dekompozycję Choleskiego dla macierzy diagonalnej:
\[
  H = \begin{bmatrix} 1 & 0 \\ 0 & 9 \end{bmatrix}
\]
otrzymujemy:
\[
  L = \begin{bmatrix} \sqrt{1} & 0 \\ 0 & \sqrt{9} \end{bmatrix} = \begin{bmatrix} 1 & 0 \\ 0 & 3 \end{bmatrix}
\]

Przemnożenie daje poprawny wynik:
\[
  L L^T = \begin{bmatrix} 1 & 0 \\ 0 & 3 \end{bmatrix} \begin{bmatrix} 1 & 0 \\ 0 & 3 \end{bmatrix} = \begin{bmatrix} 1 & 0 \\ 0 & 9 \end{bmatrix} = H
\]

\subsection{Zadanie 1(d): Minimalizacja z preconditioningiem}

\textbf{Treść zadania:} Definiując zmienne $x', y'$:
\[
  \begin{bmatrix} x' \\ y' \end{bmatrix} = L^T \begin{bmatrix} x \\ y \end{bmatrix}
\]
otrzymujemy
\[
  \begin{bmatrix} x \\ y \end{bmatrix} = L^{-T} \begin{bmatrix} x' \\ y' \end{bmatrix}
\]
oraz
\[
  g(x', y') = f\left( L^{-T} \begin{bmatrix} x' \\ y' \end{bmatrix} \right)
\]
Znajdź minimum funkcji $g$ metodą gradientu prostego. Na podstawie otrzymanego wyniku wyznacz minimum funkcji $f$.

Zastosowana transformacja preconditioning ma na celu poprawę uwarunkowania problemu. Obliczmy macierz $L^{-T}$ (równą $L^{-1}$ ponieważ $L$ jest diagonalna i symetryczna):
\[
  L^{-T} = \begin{bmatrix} 1 & 0 \\ 0 & \frac{1}{3} \end{bmatrix}
\]

Nowe współrzędne to $x' = x$ oraz $y' = 3y$. Wstawiając to do funkcji $f$, otrzymujemy nową funkcję $g(x', y')$:
\[
  g(x', y') = f(x', \tfrac{1}{3}y') = \frac{1}{2}(x')^2 + \frac{9}{2}\left( \frac{1}{3}y' \right)^2 = \frac{1}{2}(x')^2 + \frac{1}{2}(y')^2
\]

Hesjan nowej funkcji $g$ wynosi:
\[
  \nabla^2 g(x', y') = I = \begin{bmatrix} 1 & 0 \\ 0 & 1 \end{bmatrix}
\]
Współczynnik uwarunkowania tej macierzy wynosi $\kappa = 1{,}0$, co oznacza idealnie kołowe linie poziomicowe. Teoretyczny mnożnik zbieżności wynosi $R = \left( \frac{1-1}{1+1} \right)^2 = 0$. Wynika stąd, że gradient w każdym punkcie wskazuje dokładnie na minimum i algorytm powinien zbiec w dokładnie jednej iteracji.

Punkt startowy w nowej przestrzeni to:
\[
  \begin{bmatrix} x'_0 \\ y'_0 \end{bmatrix} = L^T \begin{bmatrix} x_0 \\ y_0 \end{bmatrix} = \begin{bmatrix} 1 & 0 \\ 0 & 3 \end{bmatrix} \begin{bmatrix} 9 \\ 1 \end{bmatrix} = \begin{bmatrix} 9 \\ 3 \end{bmatrix}
\]

Uruchomienie gradientu prostego z punktu $(9, 3)$ potwierdziło to, zbieżność nastąpiła po dokładnie 1 iteracji z minimum w $(0, 0)$.

\begin{lstlisting}[caption={Skalowanie zmiennych (preconditioning)}]
# L^{-T}
L_inv_T = np.linalg.inv(L.T)

def g_precond(xp_yp):
    xy = L_inv_T @ np.array(xp_yp)
    return f1(xy)

def grad_g_precond(xp_yp):
    xy = L_inv_T @ np.array(xp_yp)
    gf = grad_f1(xy)
    return np.linalg.inv(L) @ gf
\end{lstlisting}

Przeliczenie z powrotem na oryginalne zmienne za pomocą zależności $\begin{bmatrix} x \\ y \end{bmatrix} = L^{-T} \begin{bmatrix} x' \\ y' \end{bmatrix}$ daje współrzędne minimum pierwotnej funkcji $f$:
\[
  \begin{bmatrix} x \\ y \end{bmatrix} = \begin{bmatrix} 1 & 0 \\ 0 & \frac{1}{3} \end{bmatrix} \begin{bmatrix} 0 \\ 0 \end{bmatrix} = \begin{bmatrix} 0 \\ 0 \end{bmatrix}
\]

Porównanie trajektorii z preconditioningiem i bez niego przedstawia Wykres~\ref{fig:wykres2}.

\begin{figure}[H]
  \centering
  \includegraphics[width=0.95\textwidth]{wykres2.png}
  \caption{Porównanie trajektorii optymalizacji bez skalowania zmiennych (115 iteracji) oraz po zastosowaniu preconditioningu (1 iteracja).}
  \label{fig:wykres2}
\end{figure}

% ============================================================
\section{Zadanie 2: Wiszący łańcuch}
% ============================================================

\subsection{Wyprowadzenie gradientu energii potencjalnej}

\textbf{Treść zadania:} Łańcuch reprezentujemy jako $n+1$ punktowych mas o współrzędnych $(x^{(i)}, y^{(i)})$ i masie $m$ każda. Każde kolejne dwie masy połączone są sprężyną o długości w stanie spoczynku $L$ i współczynniku sprężystości $k$. Energia potencjalna układu dana jest wzorem:
\[
  V(\mathbf{x}, \mathbf{y}) = \frac{1}{2} \sum_{i=0}^{n-1} k \left( d_i - L \right)^2 + g m \sum_{i=0}^n y^{(i)}
\]
gdzie $d_i = \sqrt{(x^{(i)} - x^{(i+1)})^2 + (y^{(i)} - y^{(i+1)})^2}$ to długość $i$-tej sprężyny. Wyznacz wyrażenie na gradient $\nabla V$ funkcji celu $V$ względem położeń $x^{(i)}$ i $y^{(i)}$.

Współrzędne końców łańcucha są ustalone: $(x^{(0)}, y^{(0)}) = (0, 0)$ oraz $(x^{(n)}, y^{(n)}) = (3, 1)$. Zmiennymi decyzyjnymi w optymalizacji są wyłącznie współrzędne punktów wewnętrznych, tzn. dla $i \in \{1, \dots, n-1\}$.

Zdefiniujmy długość $i$-tej sprężyny (łączącej masę $i$ z masą $i+1$) jako:
\[
  d_i = \sqrt{(x^{(i)} - x^{(i+1)})^2 + (y^{(i)} - y^{(i+1)})^2}
\]

Pochodne cząstkowe energii sprężystości pojedynczego segmentu $V_{\text{sprz}, i} = \frac{1}{2}k(d_i - L)^2$ względem współrzędnych wynoszą:
\[
  \frac{\partial V_{\text{sprz}, i}}{\partial x^{(i)}} = k(d_i - L) \frac{x^{(i)} - x^{(i+1)}}{d_i}, \qquad \frac{\partial V_{\text{sprz}, i}}{\partial x^{(i+1)}} = -k(d_i - L) \frac{x^{(i)} - x^{(i+1)}}{d_i}
\]

Ponieważ każdy punkt wewnętrzny $i$ jest połączony z dwoma sprężynami (o indeksach $i-1$ oraz $i$), jego przesunięcie wpływa na energię obu tych sprężyn. Dodając wpływ grawitacji na współrzędną $y^{(i)}$, otrzymujemy ostateczne wzory na składowe gradientu względem położeń punktu $i$:
\[
  \frac{\partial V}{\partial x^{(i)}} = k(d_i - L) \frac{x^{(i)} - x^{(i+1)}}{d_i} - k(d_{i-1} - L) \frac{x^{(i-1)} - x^{(i)}}{d_{i-1}}
\]
\[
  \frac{\partial V}{\partial y^{(i)}} = k(d_i - L) \frac{y^{(i)} - y^{(i+1)}}{d_i} - k(d_{i-1} - L) \frac{y^{(i-1)} - y^{(i)}}{d_{i-1}} + m g
\]
Wzory te zachodzą dla wszystkich mas wewnętrznych $i = 1, \dots, n-1$.

\subsection{Zadanie 2(a): Gradient prosty ze stałym krokiem}

\textbf{Treść zadania:} Minimalizuj całkowitą energię potencjalną łańcucha metodą gradientu prostego ze stałym współczynnikiem uczenia $\alpha$. Przyjmij wartości: $n = 40$, $L = 0{,}1$ m, $m = 0{,}1$ kg, $k = 70$ N/m, $g = 9{,}81$ m/s$^2$.

Liczba stopni swobody dla $N=40$ wynosi 78 (39 wolnych punktów po 2 współrzędne każdy). Energia potencjalna dla początkowego ułożenia łańcucha (jako prostej łączącej punkty końcowe) wynosi $20{,}7246$ J.

Do optymalizacji wybrano krok $\alpha = 0{,}0005$. Dzięki poprawnej implementacji gradientu algorytm działa w pełni stabilnie i nie ulega rozbieżności nawet bez stosowania clippingu gradientu. Końcowa energia po 20000 iteracji wyniosła $-51{,}4276$ J, choć ze względu na stały krok algorytm nie osiągnął pełnej zbieżności z tolerancją $10^{-6}$ nałożoną na normę wektora gradientu.

\subsection{Zadanie 2(b): Gradient prosty z malejącym krokiem}

\textbf{Treść zadania:} Minimalizuj całkowitą energię potencjalną łańcucha metodą gradientu prostego ze współczynnikiem uczenia $\alpha$ malejącym wykładniczo co pewną liczbę iteracji.

W tej wersji algorytmu zastosowano krok początkowy $\alpha_0 = 10^{-2}$, który był zmniejszany o czynnik $0{,}7$ co 200 iteracji. Algorytm zachował pełną stabilność bez stosowania clippingu gradientu, osiągając po 20000 iteracji energię $-50{,}9222$ J. Spadek długości kroku w późniejszej fazie spowolnił dalszą optymalizację, co uniemożliwiło osiągnięcie pełnej zbieżności w ramach limitu iteracji.

\subsection{Zadanie 2(c): Gradient prosty z przeszukiwaniem liniowym}

\textbf{Treść zadania:} Minimalizuj całkowitą energię potencjalną łańcucha metodą gradientu prostego z przeszukiwaniem liniowym. Stwórz wykres ilustrujący położenie mas łańcucha po przyjęciu minimalnej energii.

W tym podejściu długość kroku $\alpha_k$ była dobierana dynamicznie w każdej iteracji przy użyciu metody backtracking line search z warunkiem Armijo. Parametry wyszukiwania to: $\alpha_0 = 1{,}0$, współczynnik zmniejszania kroku $\rho = 0{,}5$ oraz stała warunku Armijo $c = 10^{-4}$.

Dzięki dynamicznemu doborowi kroku algorytm stabilnie i szybko schodził w kierunku minimum, osiągając zbieżność (norma gradientu poniżej $10^{-6}$) po 2182 iteracjach. Końcowa energia wyniosła $-51{,}4583$ J.

Porównanie zbieżności energii oraz ostateczny kształt łańcucha przedstawiono na Wykresie~\ref{fig:wykres3}.

\begin{figure}[H]
  \centering
  \includegraphics[width=0.95\textwidth]{wykres3.png}
  \caption{Położenie mas łańcucha po optymalizacji (po lewej) oraz wykres zbieżności energii w skali logarytmicznej (po prawej) dla trzech badanych wariantów.}
  \label{fig:wykres3}
\end{figure}

Podsumowanie porównania metod optymalizacji zebrano w Tabeli~\ref{tab:tabela1}.

\begin{table}[H]
  \centering
  \caption{Porównanie metod gradientowych dla zadania wiszącego łańcucha.}
  \label{tab:tabela1}
  \begin{tabular}{lrr}
    \toprule
    Wariant algorytmu & Końcowa energia $E_{\text{final}}$ [J] & Liczba iteracji \\
    \midrule
    Stały krok ($\alpha = 0{,}0005$) & $-51{,}4276$ & 20000 (brak zbieżności) \\
    Malejący krok ($\alpha_0 = 10^{-2}$) & $-50{,}9222$ & 20000 (brak zbieżności) \\
    Przeszukiwanie liniowe (Armijo) & $-51{,}4583$ & 2182 (zbieżność) \\
    \bottomrule
  \end{tabular}
\end{table}


% ============================================================
\section{Zadanie 3: Predykcja roku wydania utworu}
% ============================================================

\textbf{Treść zadania:} Rozwiąż ponownie problem predykcji roku wydania utworu (laboratorium 2), używając metody gradientu prostego. Stałą uczącą wyznacz na podstawie najmniejszej i największej wartości własnej macierzy $A^T A$ reprezentującej cały zbiór treningowy. Porównaj uzyskane rozwiązanie z metodą najmniejszych kwadratów, biorąc pod uwagę dokładność predykcji na zbiorze testowym, teoretyczną złożoność obliczeniową oraz czas obliczeń.

Zbiór treningowy składa się z $n = 463\,715$ próbek o liczbie cech $d = 91$ (po dodaniu kolumny jedynek). Cechy oraz rok wydania zostały znormalizowane do przedziału $[0, 1]$.

\subsection{Dobór kroku uczenia}
Hesjan dla metody najmniejszych kwadratów wynosi $A^T A$. Na podstawie obliczeń numerycznych wyznaczono ekstremalne wartości własne tej macierzy:
\[
  \lambda_{\min} = 4{,}5766, \qquad \lambda_{\max} = 9{,}8235 \cdot 10^6
\]

Współczynnik uwarunkowania wynosi $\kappa \approx 2{,}15 \cdot 10^6$. Stałą uczącą $\alpha$ (zapewniającą teoretycznie najszybszą zbieżność) wyznacza się ze wzoru:
\[
  \alpha = \frac{2}{\lambda_{\min} + \lambda_{\max}} \approx 2{,}0359 \cdot 10^{-7}
\]
Niezwykle mała wartość kroku wynika bezpośrednio z bardzo złego uwarunkowania macierzy cech.

\subsection{Porównanie wyników}
Wyniki działania obu metod zebrano w Tabeli~\ref{tab:tabela2}.

\begin{table}[H]
  \centering
  \caption{Porównanie metody najmniejszych kwadratów (dokładne rozwiązanie) oraz metody gradientu prostego (po 500 iteracjach).}
  \label{tab:tabela2}
  \begin{tabular}{lrr}
    \toprule
    Kryterium & Metoda LS & Metoda GD \\
    \midrule
    \textbf{RMSE na zbiorze testowym} & $9{,}5102$ lat & $77{,}1863$ lat \\
    \textbf{Czas obliczeń} & $1{,}10$ s & $7{,}29$ s \\
    \textbf{Złożoność teoretyczna} & $\mathcal{O}(n d^2 + d^3)$ & $\mathcal{O}(M \cdot n d)$ (dla $M$ iteracji) \\
    \bottomrule
  \end{tabular}
\end{table}

\begin{itemize}
  \item \textbf{Dokładność predykcji:} Metoda najmniejszych kwadratów daje błąd testowy na poziomie $9{,}51$ roku. Z kolei metoda gradientu prostego po 500 iteracjach osiąga zaledwie $77{,}19$ roku (błąd trenujący zmniejszył się w tym czasie jedynie z $77{,}17$ do $77{,}09$ lat). Wynika to z faktu, że przy tak małym kroku $\alpha$ spadek błędu jest skrajnie wolny, do osiągnięcia zbliżonej dokładności potrzeba by milionów iteracji.
  \item \textbf{Czas obliczeń:} Metoda najmniejszych kwadratów wykonała się w $1{,}10$ s, natomiast 500 iteracji GD trwało $7{,}29$ s. GD jest wolniejszy, ponieważ w każdej iteracji musi przemnożyć duże macierze o rozmiarze $463\,715 \times 91$.
  \item \textbf{Złożoność obliczeniowa:} Teoretycznie dla bardzo dużych $d$ gradient prosty mógłby być korzystniejszy, ponieważ jego złożoność rośnie liniowo z $d$, podczas gdy LS wymaga obliczenia $A^T A$ (koszt $\mathcal{O}(n d^2)$) oraz odwrócenia tej macierzy (koszt $\mathcal{O}(d^3)$). Jednak w tym przypadku $d = 91$ jest małe, więc narzut na operacje macierzowe w LS jest znikomy, a powolna zbieżność GD eliminuje go z praktycznego zastosowania.
\end{itemize}

Wykres zbieżności błędu RMSE dla metody gradientowej w porównaniu do linii odniesienia LS przedstawiono na Wykresie~\ref{fig:wykres4}.

\begin{figure}[H]
  \centering
  \includegraphics[width=0.85\textwidth]{wykres4.png}
  \caption{Porównanie zbieżności błędu RMSE w kolejnych iteracjach metody gradientu prostego z ostatecznym wynikiem metody najmniejszych kwadratów.}
  \label{fig:wykres4}
\end{figure}

% ============================================================
\section{Wnioski}
% ============================================================

\begin{enumerate}
  \item Uwarunkowanie Hesjana ma kluczowe znaczenie dla szybkości zbieżności metody gradientu prostego. Zastosowanie dekompozycji Choleskiego jako techniki preconditioningu pozwoliło zmniejszyć liczbę iteracji z 115 do zaledwie 1, sprowadzając problem do idealnie uwarunkowanego.
  
  \item W problemach o nieliniowej strukturze i wielu stopniach swobody, takich jak wiszący łańcuch, stały krok uczenia bardzo łatwo prowadzi do rozbieżności i eksplozji energii. Konieczne jest stosowanie dynamicznych metod doboru kroku, np. przeszukiwania liniowego Armijo, które gwarantują stabilność i zbieżność do fizycznego minimum.
  
  \item W przypadku klasycznej regresji liniowej na dużych zbiorach danych, gdy liczba cech $d$ jest stosunkowo mała, bezpośrednie rozwiązanie układu normalnego za pomocą metody najmniejszych kwadratów jest nieporównywalnie szybsze i dokładniejsze niż gradient prosty. Bardzo złe uwarunkowanie macierzy cech wymusza stosowanie ekstremalnie małego kroku uczenia, przez co gradient prosty staje się bezużyteczny.
\end{enumerate}

% ============================================================
\begin{thebibliography}{2}
  \bibitem{wikipedia_gd}
  Wikipedia, \emph{Gradient descent}, \url{https://en.wikipedia.org/wiki/Gradient_descent} (dostęp: 10 czerwca 2026 r.).
  \bibitem{shewchuk_cg}
  J. R. Shewchuk, \emph{An Introduction to the Conjugate Gradient Method Without the Agonizing Pain}, Carnegie Mellon University, 1994.
\end{thebibliography}

\end{document}
